% Options for packages loaded elsewhere
\PassOptionsToPackage{unicode}{hyperref}
\PassOptionsToPackage{hyphens}{url}
\PassOptionsToPackage{dvipsnames,svgnames,x11names}{xcolor}
%
\documentclass[
  letterpaper,
  DIV=11,
  numbers=noendperiod,
  oneside]{scrartcl}

\usepackage{amsmath,amssymb}
\usepackage{iftex}
\ifPDFTeX
  \usepackage[T1]{fontenc}
  \usepackage[utf8]{inputenc}
  \usepackage{textcomp} % provide euro and other symbols
\else % if luatex or xetex
  \usepackage{unicode-math}
  \defaultfontfeatures{Scale=MatchLowercase}
  \defaultfontfeatures[\rmfamily]{Ligatures=TeX,Scale=1}
\fi
\usepackage{lmodern}
\ifPDFTeX\else  
    % xetex/luatex font selection
\fi
% Use upquote if available, for straight quotes in verbatim environments
\IfFileExists{upquote.sty}{\usepackage{upquote}}{}
\IfFileExists{microtype.sty}{% use microtype if available
  \usepackage[]{microtype}
  \UseMicrotypeSet[protrusion]{basicmath} % disable protrusion for tt fonts
}{}
\makeatletter
\@ifundefined{KOMAClassName}{% if non-KOMA class
  \IfFileExists{parskip.sty}{%
    \usepackage{parskip}
  }{% else
    \setlength{\parindent}{0pt}
    \setlength{\parskip}{6pt plus 2pt minus 1pt}}
}{% if KOMA class
  \KOMAoptions{parskip=half}}
\makeatother
\usepackage{xcolor}
\usepackage[left=1in,marginparwidth=2.0666666666667in,textwidth=4.1333333333333in,marginparsep=0.3in]{geometry}
\ifLuaTeX
  \usepackage{luacolor}
  \usepackage[soul]{lua-ul}
\else
  \usepackage{soul}
  
\fi
\setlength{\emergencystretch}{3em} % prevent overfull lines
\setcounter{secnumdepth}{-\maxdimen} % remove section numbering
% Make \paragraph and \subparagraph free-standing
\makeatletter
\ifx\paragraph\undefined\else
  \let\oldparagraph\paragraph
  \renewcommand{\paragraph}{
    \@ifstar
      \xxxParagraphStar
      \xxxParagraphNoStar
  }
  \newcommand{\xxxParagraphStar}[1]{\oldparagraph*{#1}\mbox{}}
  \newcommand{\xxxParagraphNoStar}[1]{\oldparagraph{#1}\mbox{}}
\fi
\ifx\subparagraph\undefined\else
  \let\oldsubparagraph\subparagraph
  \renewcommand{\subparagraph}{
    \@ifstar
      \xxxSubParagraphStar
      \xxxSubParagraphNoStar
  }
  \newcommand{\xxxSubParagraphStar}[1]{\oldsubparagraph*{#1}\mbox{}}
  \newcommand{\xxxSubParagraphNoStar}[1]{\oldsubparagraph{#1}\mbox{}}
\fi
\makeatother


\providecommand{\tightlist}{%
  \setlength{\itemsep}{0pt}\setlength{\parskip}{0pt}}\usepackage{longtable,booktabs,array}
\usepackage{calc} % for calculating minipage widths
% Correct order of tables after \paragraph or \subparagraph
\usepackage{etoolbox}
\makeatletter
\patchcmd\longtable{\par}{\if@noskipsec\mbox{}\fi\par}{}{}
\makeatother
% Allow footnotes in longtable head/foot
\IfFileExists{footnotehyper.sty}{\usepackage{footnotehyper}}{\usepackage{footnote}}
\makesavenoteenv{longtable}
\usepackage{graphicx}
\makeatletter
\newsavebox\pandoc@box
\newcommand*\pandocbounded[1]{% scales image to fit in text height/width
  \sbox\pandoc@box{#1}%
  \Gscale@div\@tempa{\textheight}{\dimexpr\ht\pandoc@box+\dp\pandoc@box\relax}%
  \Gscale@div\@tempb{\linewidth}{\wd\pandoc@box}%
  \ifdim\@tempb\p@<\@tempa\p@\let\@tempa\@tempb\fi% select the smaller of both
  \ifdim\@tempa\p@<\p@\scalebox{\@tempa}{\usebox\pandoc@box}%
  \else\usebox{\pandoc@box}%
  \fi%
}
% Set default figure placement to htbp
\def\fps@figure{htbp}
\makeatother

\usepackage{booktabs}
\usepackage{caption}
\usepackage{longtable}
\usepackage{colortbl}
\usepackage{array}
\usepackage{anyfontsize}
\usepackage{multirow}
\KOMAoption{captions}{tableheading}
\makeatletter
\@ifpackageloaded{caption}{}{\usepackage{caption}}
\AtBeginDocument{%
\ifdefined\contentsname
  \renewcommand*\contentsname{Table of contents}
\else
  \newcommand\contentsname{Table of contents}
\fi
\ifdefined\listfigurename
  \renewcommand*\listfigurename{List of Figures}
\else
  \newcommand\listfigurename{List of Figures}
\fi
\ifdefined\listtablename
  \renewcommand*\listtablename{List of Tables}
\else
  \newcommand\listtablename{List of Tables}
\fi
\ifdefined\figurename
  \renewcommand*\figurename{Figure}
\else
  \newcommand\figurename{Figure}
\fi
\ifdefined\tablename
  \renewcommand*\tablename{Table}
\else
  \newcommand\tablename{Table}
\fi
}
\@ifpackageloaded{float}{}{\usepackage{float}}
\floatstyle{ruled}
\@ifundefined{c@chapter}{\newfloat{codelisting}{h}{lop}}{\newfloat{codelisting}{h}{lop}[chapter]}
\floatname{codelisting}{Listing}
\newcommand*\listoflistings{\listof{codelisting}{List of Listings}}
\makeatother
\makeatletter
\makeatother
\makeatletter
\@ifpackageloaded{caption}{}{\usepackage{caption}}
\@ifpackageloaded{subcaption}{}{\usepackage{subcaption}}
\makeatother
\makeatletter
\@ifpackageloaded{sidenotes}{}{\usepackage{sidenotes}}
\@ifpackageloaded{marginnote}{}{\usepackage{marginnote}}
\makeatother

\usepackage{bookmark}

\IfFileExists{xurl.sty}{\usepackage{xurl}}{} % add URL line breaks if available
\urlstyle{same} % disable monospaced font for URLs
\hypersetup{
  pdftitle={Homework 7},
  colorlinks=true,
  linkcolor={blue},
  filecolor={Maroon},
  citecolor={Blue},
  urlcolor={Blue},
  pdfcreator={LaTeX via pandoc}}


\title{Homework 7}
\author{}
\date{}

\begin{document}
\maketitle


\subsection{Instructions}\label{instructions}

\begin{enumerate}
\def\labelenumi{(\arabic{enumi})}
\tightlist
\item
  Download the printable version of this assignment
  \href{pdfs/homework7.pdf}{here}. Print it out and record your answers
  directly in the space provided.
\item
  Write your \textbf{name} and \textbf{student number} on the top of the
  assignment.
\item
  For problems requiring additional written work, use the allotted work
  space and/or margins of the page.
\item
  Once you're done with a problem, reflect on how you well-equipped you
  felt answering that particular problem using the confidence level
  assessment.\\
\end{enumerate}

\includegraphics[width=0.75\linewidth,height=\textheight,keepaspectratio]{../images/confidence_level_assessment1.png}
\includegraphics[width=0.75\linewidth,height=\textheight,keepaspectratio]{../images/confidence_level_assessment2.png}
\includegraphics[width=0.75\linewidth,height=\textheight,keepaspectratio]{../images/confidence_level_assessment3.png}

\begin{enumerate}
\def\labelenumi{(\arabic{enumi})}
\setcounter{enumi}{4}
\tightlist
\item
  Turn in your assignment in class on the due date.
\end{enumerate}

\subsubsection{Questions}\label{questions}

\paragraph{Question 1}\label{question-1}

\textbf{A researcher collected the following data on years of education
(X) and number of children (Y) for a sample of married adults:}

\begin{longtable}[]{@{}
  >{\centering\arraybackslash}p{(\linewidth - 2\tabcolsep) * \real{0.5000}}
  >{\centering\arraybackslash}p{(\linewidth - 2\tabcolsep) * \real{0.5000}}@{}}
\toprule\noalign{}
\begin{minipage}[b]{\linewidth}\centering
X
\end{minipage} & \begin{minipage}[b]{\linewidth}\centering
Y
\end{minipage} \\
\midrule\noalign{}
\endhead
\bottomrule\noalign{}
\endlastfoot
12 & 2 \\
14 & 1 \\
17 & 0 \\
10 & 3 \\
8 & 5 \\
9 & 3 \\
12 & 4 \\
14 & 2 \\
18 & 0 \\
16 & 2 \\
\end{longtable}

\begin{enumerate}
\def\labelenumi{\alph{enumi}.}
\tightlist
\item
  Draw a scatter plot of the data.
\end{enumerate}

\marginnote{\begin{footnotesize}

\includegraphics[width=0.75\linewidth,height=\textheight,keepaspectratio]{../images/confidence_level_assessment.png}

\end{footnotesize}}

\begin{enumerate}
\def\labelenumi{\alph{enumi}.}
\setcounter{enumi}{1}
\tightlist
\item
  Calculate and interpret the correlation between education and number
  of children.
\end{enumerate}

\marginnote{\begin{footnotesize}

\includegraphics[width=0.75\linewidth,height=\textheight,keepaspectratio]{../images/confidence_level_assessment.png}

\end{footnotesize}}

\begin{enumerate}
\def\labelenumi{\alph{enumi}.}
\setcounter{enumi}{2}
\tightlist
\item
  Is the association between education and number of children
  statistically significant? Explain the basis of your answer and the
  implications for the association in the population.
\end{enumerate}

\marginnote{\begin{footnotesize}

\includegraphics[width=0.75\linewidth,height=\textheight,keepaspectratio]{../images/confidence_level_assessment.png}

\end{footnotesize}}

\begin{enumerate}
\def\labelenumi{\alph{enumi}.}
\setcounter{enumi}{3}
\tightlist
\item
  Calculate and interpret the meaning of the regression slope.
\end{enumerate}

\marginnote{\begin{footnotesize}

\includegraphics[width=0.75\linewidth,height=\textheight,keepaspectratio]{../images/confidence_level_assessment.png}

\end{footnotesize}}

\begin{enumerate}
\def\labelenumi{\alph{enumi}.}
\setcounter{enumi}{4}
\tightlist
\item
  Calculate and interpret the meaning of the Y-intercept for the
  regression line.
\end{enumerate}

\marginnote{\begin{footnotesize}

\includegraphics[width=0.75\linewidth,height=\textheight,keepaspectratio]{../images/confidence_level_assessment.png}

\end{footnotesize}}

\begin{enumerate}
\def\labelenumi{\alph{enumi}.}
\setcounter{enumi}{5}
\tightlist
\item
  Predict the number of children for an adult with 11 years of
  education.
\end{enumerate}

\marginnote{\begin{footnotesize}

\includegraphics[width=0.75\linewidth,height=\textheight,keepaspectratio]{../images/confidence_level_assessment.png}

\end{footnotesize}}

\begin{enumerate}
\def\labelenumi{\alph{enumi}.}
\setcounter{enumi}{6}
\tightlist
\item
  Calculate and interpret the coefficient of determination.
\end{enumerate}

\marginnote{\begin{footnotesize}

\includegraphics[width=0.75\linewidth,height=\textheight,keepaspectratio]{../images/confidence_level_assessment.png}

\end{footnotesize}}

\paragraph{Question 2}\label{question-2}

\textbf{The same researcher collected at random one child from each of
the eight families in which there was at least one child. She calculated
the following statistics to study the effects of number of siblings (X)
on happiness of the child (Y), measured on an interval scale in which
low scores indicate a low level of happiness and high scores indicate a
high level of happiness.}

\begin{longtable}[]{@{}
  >{\centering\arraybackslash}p{(\linewidth - 4\tabcolsep) * \real{0.2000}}
  >{\centering\arraybackslash}p{(\linewidth - 4\tabcolsep) * \real{0.2000}}
  >{\centering\arraybackslash}p{(\linewidth - 4\tabcolsep) * \real{0.2000}}@{}}
\toprule\noalign{}
\begin{minipage}[b]{\linewidth}\centering
Happiness points (Y)
\end{minipage} & \begin{minipage}[b]{\linewidth}\centering
Number of siblings (X)
\end{minipage} & \begin{minipage}[b]{\linewidth}\centering
\end{minipage} \\
\midrule\noalign{}
\endhead
\bottomrule\noalign{}
\endlastfoot
\(\bar{y} = 5.625\) & \(\bar{x} = 1.750\) & \(r = 0.588\) \\
\(s_y = 2.134\) & \(s_x = 1.282\) & \(n = 8\) \\
\end{longtable}

\begin{enumerate}
\def\labelenumi{\alph{enumi}.}
\tightlist
\item
  Is the association between number of siblings and happiness
  statistically significant? Explain the basis of your answer and the
  implications for the association in the population.
\end{enumerate}

\marginnote{\begin{footnotesize}

\includegraphics[width=0.75\linewidth,height=\textheight,keepaspectratio]{../images/confidence_level_assessment.png}

\end{footnotesize}}

\begin{enumerate}
\def\labelenumi{\alph{enumi}.}
\setcounter{enumi}{1}
\tightlist
\item
  Calculate and interpret the meaning of the regression slope.
\end{enumerate}

\marginnote{\begin{footnotesize}

\includegraphics[width=0.75\linewidth,height=\textheight,keepaspectratio]{../images/confidence_level_assessment.png}

\end{footnotesize}}

\begin{enumerate}
\def\labelenumi{\alph{enumi}.}
\setcounter{enumi}{2}
\tightlist
\item
  Calculate and interpret the meaning of the Y-intercept for the
  regression line.
\end{enumerate}

\marginnote{\begin{footnotesize}

\includegraphics[width=0.75\linewidth,height=\textheight,keepaspectratio]{../images/confidence_level_assessment.png}

\end{footnotesize}}

\begin{enumerate}
\def\labelenumi{\alph{enumi}.}
\setcounter{enumi}{3}
\tightlist
\item
  Predict the level of happiness for an only child and for a kid with
  two siblings.
\end{enumerate}

\marginnote{\begin{footnotesize}

\includegraphics[width=0.75\linewidth,height=\textheight,keepaspectratio]{../images/confidence_level_assessment.png}

\end{footnotesize}}

\begin{enumerate}
\def\labelenumi{\alph{enumi}.}
\setcounter{enumi}{4}
\tightlist
\item
  Calculate and interpret the coefficient of determination.
\end{enumerate}

\marginnote{\begin{footnotesize}

\includegraphics[width=0.75\linewidth,height=\textheight,keepaspectratio]{../images/confidence_level_assessment.png}

\end{footnotesize}}

\paragraph{Question 3}\label{question-3}

\textbf{A personnel specialist with a large accounting firm is
interested in determining the effect of seniority (the number of years
with the company) on hourly wages for administrative staff. She selects
a random sample of 10 secretaries and collects the following data on
their years with the company (X) and hourly wages (Y):}

\begin{longtable}[]{@{}
  >{\centering\arraybackslash}p{(\linewidth - 4\tabcolsep) * \real{0.2000}}
  >{\centering\arraybackslash}p{(\linewidth - 4\tabcolsep) * \real{0.2000}}
  >{\centering\arraybackslash}p{(\linewidth - 4\tabcolsep) * \real{0.2000}}@{}}
\toprule\noalign{}
\begin{minipage}[b]{\linewidth}\centering
Hourly Wage (Y)
\end{minipage} & \begin{minipage}[b]{\linewidth}\centering
Years with company (X)
\end{minipage} & \begin{minipage}[b]{\linewidth}\centering
\end{minipage} \\
\midrule\noalign{}
\endhead
\bottomrule\noalign{}
\endlastfoot
\(\bar{y} = \$13.90\) & \(\bar{x} = 2.700\) & \(r = 0.631\) \\
\(s_y = \$1.37\) & \(s_x = 1.889\) & \(n = 10\) \\
\end{longtable}

\begin{enumerate}
\def\labelenumi{\alph{enumi}.}
\tightlist
\item
  Is the association statistically significant? Explain your answer.
\end{enumerate}

\marginnote{\begin{footnotesize}

\includegraphics[width=0.75\linewidth,height=\textheight,keepaspectratio]{../images/confidence_level_assessment.png}

\end{footnotesize}}

\begin{enumerate}
\def\labelenumi{\alph{enumi}.}
\setcounter{enumi}{1}
\tightlist
\item
  Calculate and interpret the meaning of the regression slope.
\end{enumerate}

\marginnote{\begin{footnotesize}

\includegraphics[width=0.75\linewidth,height=\textheight,keepaspectratio]{../images/confidence_level_assessment.png}

\end{footnotesize}}

\begin{enumerate}
\def\labelenumi{\alph{enumi}.}
\setcounter{enumi}{2}
\tightlist
\item
  Calculate and interpret the meaning of the Y-intercept for the
  regression line.
\end{enumerate}

\marginnote{\begin{footnotesize}

\includegraphics[width=0.75\linewidth,height=\textheight,keepaspectratio]{../images/confidence_level_assessment.png}

\end{footnotesize}}

\begin{enumerate}
\def\labelenumi{\alph{enumi}.}
\setcounter{enumi}{3}
\tightlist
\item
  Predict the hourly wage of a randomly selected worker who has been
  with the company for 4 years.
\end{enumerate}

\marginnote{\begin{footnotesize}

\includegraphics[width=0.75\linewidth,height=\textheight,keepaspectratio]{../images/confidence_level_assessment.png}

\end{footnotesize}}

\begin{enumerate}
\def\labelenumi{\alph{enumi}.}
\setcounter{enumi}{4}
\tightlist
\item
  Calculate and interpret the coefficient of determination.
\end{enumerate}

\marginnote{\begin{footnotesize}

\includegraphics[width=0.75\linewidth,height=\textheight,keepaspectratio]{../images/confidence_level_assessment.png}

\end{footnotesize}}

\begin{enumerate}
\def\labelenumi{\alph{enumi}.}
\setcounter{enumi}{5}
\tightlist
\item
  Based on these results, what is the typical starting wage per hour,
  and what is the typical increase in wage for each additional year on
  the job? (\emph{not a trick question!})
\end{enumerate}

\marginnote{\begin{footnotesize}

\includegraphics[width=0.75\linewidth,height=\textheight,keepaspectratio]{../images/confidence_level_assessment.png}

\end{footnotesize}}

\paragraph{Use the following information to answer both questions 4 and
5}\label{use-the-following-information-to-answer-both-questions-4-and-5}

\textbf{Even in the modern era, countries differ dramatically in terms
of the health of their populations. For example, average life expectancy
is in the 80's in some countries but remains in the low 40's in other
countries. Some researchers claim that the most effective way to improve
the health of the worst-off populations is to provide more education for
people in these countries. To assess this argument, use the following
analysis of data from a random sample of countries}

\begin{figure*}

\begin{table}
\caption*{
{\large Results of OLS regression models predicting<br>life expectancy at birth in a sample of countries}
} 
\fontsize{12.0pt}{14.4pt}\selectfont
\begin{tabular*}{1\linewidth}{@{\extracolsep{\fill}}l|rr}
\toprule
 & Model 1 & Model 2 \\ 
\midrule\addlinespace[2.5pt]
\multicolumn{3}{>{\raggedright\arraybackslash}m{1\linewidth}}{Predictor} \\[2.5pt] 
\midrule\addlinespace[2.5pt]
{\hspace{15pt}\% Literate} & 0.403 *** & 0.301 *** \\ 
{\hspace{15pt}\% Urban} &  & 0.149 *** \\ 
\midrule\addlinespace[2.5pt]
\multicolumn{3}{>{\raggedright\arraybackslash}m{1\linewidth}}{\rule{0pt}{0pt}} \\[-3.2ex] 
\midrule\addlinespace[2.5pt]
Constant & 38.541 & 37.974 \\ 
Model R² & 0.749 & 0.814 \\ 
\bottomrule
\end{tabular*}
\begin{minipage}{\linewidth}
N = 107 <br><br> * p < 0.05 &ensp; ** p < 0.01 &ensp; *** p < 0.001\\
\end{minipage}
\end{table}

\ul{\textbf{Definition of Variables}}

\begin{itemize}
\tightlist
\item
  Life expectancy at birth: Average age to which people born right now
  will live if exposed to current death rates in the country.
\item
  \% Literate: Percent of adults in the population that can read and
  write.
\item
  \% Urban: Percent of the country's population residing in cities.
\end{itemize}

\end{figure*}%

\newpage{}

\paragraph{Question 4}\label{question-4}

\textbf{Interpretation of Model 1}

\begin{enumerate}
\def\labelenumi{\alph{enumi}.}
\tightlist
\item
  Write out the regression equation implied in the results for Model 1.
\end{enumerate}

\marginnote{\begin{footnotesize}

\includegraphics[width=0.75\linewidth,height=\textheight,keepaspectratio]{../images/confidence_level_assessment.png}

\end{footnotesize}}

\begin{enumerate}
\def\labelenumi{\alph{enumi}.}
\setcounter{enumi}{1}
\tightlist
\item
  Predict the life expectancy for a country in which 75\% of the
  population can read and write.
\end{enumerate}

\marginnote{\begin{footnotesize}

\includegraphics[width=0.75\linewidth,height=\textheight,keepaspectratio]{../images/confidence_level_assessment.png}

\end{footnotesize}}

\begin{enumerate}
\def\labelenumi{\alph{enumi}.}
\setcounter{enumi}{2}
\tightlist
\item
  Interpret the meaning of the constant in Model 1.
\end{enumerate}

\marginnote{\begin{footnotesize}

\includegraphics[width=0.75\linewidth,height=\textheight,keepaspectratio]{../images/confidence_level_assessment.png}

\end{footnotesize}}

\begin{enumerate}
\def\labelenumi{\alph{enumi}.}
\setcounter{enumi}{3}
\tightlist
\item
  Interpret the slope coefficient for the effect of \% Literate in Model
  1.
\end{enumerate}

\marginnote{\begin{footnotesize}

\includegraphics[width=0.75\linewidth,height=\textheight,keepaspectratio]{../images/confidence_level_assessment.png}

\end{footnotesize}}

\begin{enumerate}
\def\labelenumi{\alph{enumi}.}
\setcounter{enumi}{4}
\tightlist
\item
  Is the association between literacy and life expectancy statistically
  significant? Explain how you reached your conclusion.
\end{enumerate}

\marginnote{\begin{footnotesize}

\includegraphics[width=0.75\linewidth,height=\textheight,keepaspectratio]{../images/confidence_level_assessment.png}

\end{footnotesize}}

\begin{enumerate}
\def\labelenumi{\alph{enumi}.}
\setcounter{enumi}{5}
\tightlist
\item
  Interpret the coefficient of determination in Model 1.
\end{enumerate}

\marginnote{\begin{footnotesize}

\includegraphics[width=0.75\linewidth,height=\textheight,keepaspectratio]{../images/confidence_level_assessment.png}

\end{footnotesize}}

\begin{enumerate}
\def\labelenumi{\alph{enumi}.}
\setcounter{enumi}{6}
\tightlist
\item
  Do these results support the idea that population health (as measured
  by life expectancy) is affected by education levels (as measured by \%
  literate)? Explain your answer.
\end{enumerate}

\marginnote{\begin{footnotesize}

\includegraphics[width=0.75\linewidth,height=\textheight,keepaspectratio]{../images/confidence_level_assessment.png}

\end{footnotesize}}

\paragraph{Question 5}\label{question-5}

\textbf{A critic of the education argument suggests that any association
between literacy and life expectancy is a spurious byproduct of
urbanization; in countries in which the population is concentrated in
cities people have access to both schooling and health care, creating an
association between literacy between literacy and life expectancy that
is not causal. To test this argument, interpret Model 2 of the table
above.}

\begin{enumerate}
\def\labelenumi{\alph{enumi}.}
\tightlist
\item
  Write out the regression equation implied in the results for Model 2.
\end{enumerate}

\marginnote{\begin{footnotesize}

\includegraphics[width=0.75\linewidth,height=\textheight,keepaspectratio]{../images/confidence_level_assessment.png}

\end{footnotesize}}

\begin{enumerate}
\def\labelenumi{\alph{enumi}.}
\setcounter{enumi}{1}
\tightlist
\item
  Predict the life expectancy for a country in which 75\% of the
  population can read and write and 50\% of the population lives in
  cities.
\end{enumerate}

\marginnote{\begin{footnotesize}

\includegraphics[width=0.75\linewidth,height=\textheight,keepaspectratio]{../images/confidence_level_assessment.png}

\end{footnotesize}}

\begin{enumerate}
\def\labelenumi{\alph{enumi}.}
\setcounter{enumi}{2}
\tightlist
\item
  Interpret the meaning of the constant in Model 2.
\end{enumerate}

\marginnote{\begin{footnotesize}

\includegraphics[width=0.75\linewidth,height=\textheight,keepaspectratio]{../images/confidence_level_assessment.png}

\end{footnotesize}}

\begin{enumerate}
\def\labelenumi{\alph{enumi}.}
\setcounter{enumi}{3}
\tightlist
\item
  Interpret the slope coefficient for the effect of \% Literate in Model
  2.
\end{enumerate}

\marginnote{\begin{footnotesize}

\includegraphics[width=0.75\linewidth,height=\textheight,keepaspectratio]{../images/confidence_level_assessment.png}

\end{footnotesize}}

\begin{enumerate}
\def\labelenumi{\alph{enumi}.}
\setcounter{enumi}{4}
\tightlist
\item
  Interpret the slope coefficient for the effect of \% Urban in Model 2.
\end{enumerate}

\marginnote{\begin{footnotesize}

\includegraphics[width=0.75\linewidth,height=\textheight,keepaspectratio]{../images/confidence_level_assessment.png}

\end{footnotesize}}

\begin{enumerate}
\def\labelenumi{\alph{enumi}.}
\setcounter{enumi}{5}
\tightlist
\item
  Interpret the coefficient of multiple determination in Model 2
\end{enumerate}

\marginnote{\begin{footnotesize}

\includegraphics[width=0.75\linewidth,height=\textheight,keepaspectratio]{../images/confidence_level_assessment.png}

\end{footnotesize}}

\begin{enumerate}
\def\labelenumi{\alph{enumi}.}
\setcounter{enumi}{6}
\tightlist
\item
  Do these results support the argument that the association between
  life expectancy and \% literate is spurious? Explain your answer.
\end{enumerate}

\marginnote{\begin{footnotesize}

\includegraphics[width=0.75\linewidth,height=\textheight,keepaspectratio]{../images/confidence_level_assessment.png}

\end{footnotesize}}

\end{document}
